\subsection{EV Scheduling Agent}

\subsubsection*{\textbf{Problem Description}}
Day-ahead EV charging scheduling requires allocating power to concurrent sessions so that each vehicle receives its requested energy before departure while respecting per-charger power limits and a site-level capacity cap. Formally, sessions $\mathcal{I}$ are scheduled over a discrete horizon $\mathcal{T} = \{0,\ldots,95\}$ (96 steps of 15\,min, one day). Each session $i \in \mathcal{I}$ has arrival $a_i$, departure $d_i$, requested energy $E_i^{\mathrm{req}}$ [kWh], and per-session power cap $\bar{p}_i$ [kW]. The site has aggregate cap $\bar{P}_t$ [kW] and TOU price $c_t$ [\$/kWh]:
\begin{align}
\min_{\{p_{i,t}\},\{u_i\}} \quad &
    \Delta t \sum_t c_t \sum_i p_{i,t}
    + \lambda \sum_i u_i \label{eq:ev_obj}\\
\text{s.t.}\quad
& p_{i,t} = 0, \quad t \notin [a_i, d_i) \label{eq:ev_avail}\\
& 0 \le p_{i,t} \le \bar{p}_i \label{eq:ev_power}\\
& \textstyle\sum_i p_{i,t} \le \bar{P}_t \label{eq:ev_cap}\\
& \Delta t \textstyle\sum_{t=a_i}^{d_i-1} p_{i,t} + u_i
    = E_i^{\mathrm{req}}, \quad u_i \ge 0. \label{eq:ev_energy}
\end{align}
The slack $u_i \ge 0$ permits unmet energy only when the site cap makes full delivery physically impossible.
The evaluation uses 19 benchmark days from the Caltech ACN-Data JPL site~\cite{lee2019acndata}, covering 2018--2019 with a 50\,kW cap, median 20 sessions per day (range 4--47), and median 15\,kWh per session. TOU prices are \$0.12/kWh off-peak (steps 0--63, 84--95) and \$0.45/kWh peak (steps 64--83, i.e., 4--9\,pm).

\subsubsection*{\textbf{Method}}

We compare three approaches on identical session data, all using the same CVXPY solver to solve~(\ref{eq:ev_obj})--(\ref{eq:ev_energy}) with $\Delta t = 0.25$\,h and $\lambda = 10^6$\,\$/kWh (linear program, globally optimal)~\cite{diamond2016cvxpy}:

\begin{enumerate}[leftmargin=*,label=(\Alph*),itemsep=1pt]
\item \textit{Numerical solver}: CVXPY called directly on structured session data (gold standard).
\item \textit{LLM-only baseline}: GPT-4o receives the problem as natural-language text and must output the schedule matrix directly. No optimizer or tools accessible during generation.
\item \textit{\textsc{AgenticEV}}: GPT-4o with tool access to the CVXPY solver via OpenAI function calling. The LLM autonomously decides when to invoke the optimizer tool versus answering directly.
\end{enumerate}

All three pipelines share an independent constraint checker for post-solve validation, ensuring performance differences reflect tool access rather than evaluation bias (Fig.~\ref{fig:ev_architecture}).

% QUESTION 1 FOR STUDENTS: Your report describes a "5-step ReAct loop: Parse → Plan → Optimize →
% Validate → Refine → Explain". However, the code shows Plan (agent/plan/plan.py:34-63) and
% Refine (agent/refine/refine.py:17-43) are pass-throughs with no logic (comments say "v1:
% pass-through, no refinement logic"). Why does the report claim 5 steps when only 3 are implemented?
% Was this initial design that wasn't finished, or a misunderstanding of what was built?

\textsc{AgenticEV} operates in three stages, with stages 2--3 within a multi-turn conversation loop (up to 3 tool-call rounds):

\begin{enumerate}[leftmargin=*,label=\arabic*,itemsep=1pt]
\item \textit{Parse (GPT-4o).} A dedicated LLM call with structured output parsing converts the free-form NL request into session parameters: step-indexed arrival/departure windows, per-session energy, and per-session max power. Informal time expressions (e.g., ``half past six'') are mapped to step indices. Missing fields trigger a clarification request.

\item \textit{Optimize (CVXPY tool).} Within a multi-turn conversation, the LLM autonomously decides whether to invoke \texttt{solve\_ev\_schedule} based on the request type (scheduling query vs.\ conceptual question). When called, the tool runs the convex solver on~(\ref{eq:ev_obj})--(\ref{eq:ev_energy}) and returns a JSON summary: \texttt{total\_cost\_usd}, \texttt{peak\_load\_kw}, \texttt{total\_unmet\_kwh}, \texttt{pct\_fully\_served}. The full schedule matrix is stored internally but not exposed to the LLM. The LLM can call the tool multiple times (e.g., for what-if scenarios via \texttt{disabled\_chargers}, \texttt{site\_cap\_kw}, \texttt{extra\_sessions} parameters).

\item \textit{Explain (GPT-4o).} The LLM synthesizes the tool's JSON output into a natural-language summary. All numerical claims in the explanation are traceable to specific tool return values, structurally preventing metric hallucination.
\end{enumerate}

% QUESTION 2 FOR STUDENTS: Your report states "Validate (checker): any flags are surfaced to the
% agent" and "Refine (GPT-4o): corrects and retries". However, validate() is called on line 395
% of agent/llm_agent.py, AFTER the agent loop exits (lines 332-376). The validation results are
% NOT added to messages[], NOT seen by the LLM, and NOT used for self-correction. Why does the
% report claim validation is part of the agent loop when it's actually post-processing? Was
% validation intended to be in the loop originally? If so, what prevented implementation?

A post-hoc constraint checker independently verifies feasibility after agent completion; validation results are not fed back to the agent. This architecture enforces the solver-grounded rule: every reported number originates from the verified CVXPY tool, not from LLM inference. Fig.~\ref{fig:ev_trajectory} shows a complete \textsc{AgenticEV} trajectory for a 7-EV natural-language request.

\begin{figure}[htb]
\centering
\begin{minipage}{0.94\linewidth}
\begin{promptfigure}{AgenticEV Trajectory, 7-EV Free-Form Request (Caltech ACN-Data JPL Site)}
\footnotesize
{\ttfamily\small
\rolehl{User:} I have 7 EVs: EV1 arrives 6\,pm leaves 10\,pm needs 20\,kWh, EV2 arrives 7\,pm leaves 11\,pm needs 15\,kWh, EV3 arrives 8\,pm leaves midnight needs 10\,kWh, EV4 arrives 3\,pm leaves 4\,pm needs 3\,kWh, EV5 arrives 5\,pm leaves 7:30\,pm needs 3\,kWh, EV6 arrives noon leaves 6\,pm needs 30\,kWh, EV7 arrives 8\,am leaves 11\,am needs 6\,kWh.\par\vspace{1mm}
\reasonhl{Agent (Parse):} Mapped 7 sessions to 15-min step indices. EV1: steps 72--87, EV2: steps 76--92, EV3: steps 80--95, EV4: steps 60--63, EV5: steps 68--74, EV6: steps 48--72, EV7: steps 32--44. Per-session power caps loaded from charger metadata.\par\vspace{1mm}
\contexthl{Tool call:} \texttt{solve\_ev\_schedule(\{sessions: [7 entries], site\_cap\_kw: 50, tou\_off\_peak: 0.12, tou\_peak: 0.45\})}\par\vspace{1mm}
\retrievedhl{Tool output:} \texttt{\{total\_cost\_usd: 18.96, peak\_load\_kw: 17.38, total\_unmet\_kwh: 0.0, pct\_fully\_served: 100.0\}}\par\vspace{1mm}
\contexthl{Checker:} 0 hard constraint violations; 0 energy shortfall flags.\par\vspace{1mm}
\outputhl{Agent (Explain):} ``The charging schedule serves all 7 sessions fully with no unmet energy. Total cost: \$18.96. Peak load: 17.38\,kW, well below the 50\,kW site cap. Sessions with earlier departure windows (EV4, EV7) are charged off-peak to minimize TOU cost.''}
\end{promptfigure}
\end{minipage}
\caption{Complete \textsc{AgenticEV} trajectory for a 7-EV free-form request. The trace shows end-to-end system behavior: (1)~user input, (2)~preprocessing via \texttt{parse\_nl\_problem()} (separate GPT-4o call), (3)~agent conversation loop where the LLM invokes \texttt{solve\_ev\_schedule} via function calling, (4)~tool returns JSON metrics, (5)~agent generates grounded explanation. All numbers (18.96, 17.38, 100\%) are copied from tool output, not inferred by the LLM.}
\label{fig:ev_trajectory}
\end{figure}

\subsubsection*{\textbf{Results}}

We compare three methods on identical session data scored by the same checker. Checker flags combine two types: \textit{hard constraint violations} on~(\ref{eq:ev_avail})--(\ref{eq:ev_cap}) (physical infeasibility) and \textit{energy shortfall flags} where $u_i > 0$ (under-delivery forced by the site cap).

\begin{table}[ht]
\centering
\caption{Benchmark results over 19 days (Caltech ACN-Data JPL site).
Bold = best per column. Lower LLM-baseline cost is an artifact of
under-delivery, not genuine optimization.}
\label{tab:ev_results}
\footnotesize
\renewcommand{\arraystretch}{1.12}
\begin{tabularx}{\linewidth}{l Y Y Y Y Y}
\toprule
\textbf{Method}
  & \textbf{Avg Cost (\$)~$\downarrow$}
  & \textbf{Avg Peak (kW)~$\downarrow$}
  & \textbf{Avg Unmet (kWh)~$\downarrow$}
  & \textbf{Avg Served (\%)~$\uparrow$}
  & \textbf{Avg Flags~$\downarrow$} \\
  \midrule
  Numerical Solver   & \textbf{97.19} & \textbf{47.21}
  & \textbf{20.95} & \textbf{73.86}
  & \textbf{13.11} \\
  LLM-Only Baseline  & $84.34^\dagger$ & 60.97
  & 251.59 & 31.83 & 86.11 \\
  \textsc{AgenticEV} & \textbf{97.19} & \textbf{47.21}
  & \textbf{20.95} & \textbf{73.86}
  & \textbf{13.11} \\
  \bottomrule
  \multicolumn{6}{p{\linewidth}}{\scriptsize
  $^\dagger$Lower cost is misleading: the baseline under-delivers
  251.59\,kWh vs.\ 20.95\,kWh for the solver and agent
  ($12\times$ worse). Baseline flags include hard physical-constraint
  violations; solver/agent flags are entirely energy shortfalls on
  congested days (zero hard violations). $n = 19$ days.}
  \end{tabularx}
  \end{table}

Table~\ref{tab:ev_results} shows that \textsc{AgenticEV} exactly matches the numerical solver on all five metrics across all 19 days. The result directly validates the tool-augmented design principle: by delegating optimization to the CVXPY tool via function calling, the agent inherits the convex solver's global optimality guarantee without the LLM performing any numerical reasoning. The LLM-only baseline fails on all three success criteria—12$\times$ more unmet energy, 6.6$\times$ more checker flags, and hard physical-constraint violations including peak load exceeding the 50\,kW site cap—confirming that a standalone LLM cannot replace a solver for constrained scheduling tasks. Its apparently lower cost is entirely an artifact of under-delivery. \textsc{AgenticEV} completes at approximately 30\,s per day (3--5 calls, 2--5k tokens each) compared to $<$1\,s for the direct solver—the overhead of natural-language accessibility.

\clearpage
